\section{World-Model Quality and the Lyapunov Function}
\label{sec:lyapunov}

The phase boundary requires a scalar measure of world-model quality
whose dynamics can be bounded.  We define this as a weighted
sum-of-squares error and establish the per-step bounds that drive the
rest of the analysis.

\subsection{Measuring world-model quality}

\begin{definition}[World-Model Error]
\label{def:wm_error}
The error of $\Wm_t$ on dimension $k$ is
\begin{equation}
\label{eq:dim_error}
\epsilon_k(t) = |\Wm_t[k] - \Wm^*[k]|
\end{equation}
where $\Wm^*[k]$ is the true value of dimension $k$ (the true leverage
of capability $d$, the true risk of tuple $S$, the true effective
substrate diversity, etc.).  The total error is:
\begin{equation}
\label{eq:lyapunov}
\Lyap(\Wm_t) = \sum_{k=1}^{K} w_k \, \epsilon_k(t)^2
\end{equation}
where $w_k > 0$ is the safety-relevance weight of dimension $k$
(higher for dimensions whose miscalibration produces larger $\volL$
errors).  $\Lyap$ serves as the Lyapunov candidate.
\end{definition}

\begin{remark}[Weight choice]
\label{rem:weights}
Two natural choices for $w_k$ present themselves.  The first is
$\volL$-sensitivity: $w_k = |\partial \volL / \partial \Wm[k]|$,
making $\Lyap$ measure $\volL$-relevant error directly.  The second is
uniform weights ($w_k = 1$ for all $k$), which simplifies the analysis.
Both produce valid phase boundaries; the choice affects the constants
$c_{\mathrm{S}}$ and $c_{\mathrm{D}}$ in the comparison-function
assumptions below but not the structural form of the phase boundary
theorem.  We leave $w_k$ as a free parameter and note that the
$\volL$-sensitivity choice is the natural default for a framework whose
objective is $\volL$ maximization.
\end{remark}

\subsection{Lyapunov dynamics}

The per-step behavior of $\Lyap$ depends on two comparison-function
assumptions that bound the correction and degradation signals.

\begin{assumption}[Coercivity of Self-Correction (C1)]
\label{ass:C1}
There exists a constant $c_{\mathrm{S}} > 0$ such that the aggregate
correction signal satisfies
\begin{equation}
\label{eq:C1}
S(t) \geq c_{\mathrm{S}} \cdot \rhomincross(P_t) \cdot \Lyap(\Wm_t)
\end{equation}
whenever $\Lyap(\Wm_t) > 0$ and $\rhomincross(P_t) > 0$.
\end{assumption}

This states that aggregate correction is at least proportional to both
the current error ($\Lyap$) and the available channel strength
($\rhomincross$).  The justification follows from the EWMA learning
rule structure in Paper~4's risk-trust model: each observed channel
drives its covered dimensions toward truth at a rate proportional to
the error on that dimension (linear in the residual) and to the
channel's effective sensitivity (captured in $\rho_k$).  Additivity
across cross-substrate channels follows from Paper~5's channel
independence property under the joint-factorization assumption.

\begin{assumption}[Bounded Degradation (C2$'$)]
\label{ass:C2prime}
There exist constants $c_{\mathrm{D}} > 0$ and $\dmax > 0$ such that
for any step $t$ (subsumption or otherwise):
\begin{equation}
\label{eq:C2prime}
D(t) \leq c_{\mathrm{D}} \cdot \nsub(t) \cdot \Lyap(\Wm_t)
  + \dmax \cdot \Delta\rho(t)
\end{equation}
where $\nsub(t) \in \{0,1\}$ indicates whether a subsumption occurred
at step $t$, and
\begin{equation}
\label{eq:delta_rho}
\Delta\rho(t) = \sum_{k=1}^{K} w_k \cdot
  \max\!\left(0,\; \rhocross_k(P_t) - \rhocross_k(P_{t+1})\right)
\end{equation}
is the weighted redundancy lost at step $t$.
\end{assumption}

Assumption~C2$'$ contains two terms.  The first, $c_{\mathrm{D}} \cdot
\Lyap(\Wm_t)$, captures error-proportional growth: existing errors
compound faster without correction.  The second, $\dmax \cdot
\Delta\rho(t)$, captures the \emph{new} drift exposure from lost
channels.  This second term is the critical addition relative to the
simpler ``C2'' one might write without it.  A dimension with
$\epsilon_k \approx 0$ (well-calibrated \emph{because} a channel was
actively correcting it) starts drifting at rate $\dmax$ when its
channel is removed, even though $\Lyap(\Wm_t) \approx 0$ at the moment
of removal.  Without the $\Delta\rho$ term, the bound would predict
near-zero degradation precisely when subsumption is most
dangerous---removing a channel that was doing all the correction work
on a dimension that would otherwise drift.

\begin{proposition}[Error Dynamics]
\label{prop:error_dynamics}
Under the coupled $(P, \Wm)$ dynamics with locally rational transitions
(Definition~\ref{def:transition}), Assumptions~C1 and~C2$'$, and
$\alpha(t)$, $\beta(t)$ denoting the learning rate and degradation
coefficient respectively, the Lyapunov candidate satisfies two per-step
bounds.

\medskip
\noindent\textbf{Non-subsumption steps} ($\nsub(t) = 0$):
\begin{equation}
\label{eq:contraction}
\Lyap(\Wm_{t+1}) \leq \Lyap(\Wm_t) \cdot
  \left(1 - \rS \cdot \rhomincross(P_t)\right)
\end{equation}
where $\rS = \alpha(t) \cdot c_{\mathrm{S}}$ is the self-correction
rate.  This is a strict contraction provided $0 < \rS \cdot
\rhomincross(P_t) < 1$ (i.e., the step size is small enough that
correction does not overshoot).  We assume throughout that $\alpha(t)$
is chosen small enough to satisfy this bound; this is the standard
step-size condition for EWMA-style learning rules and holds whenever
$\alpha(t) \leq 1 / (c_{\mathrm{S}} \cdot \rhomincross(P_t))$.

\medskip
\noindent\textbf{Subsumption steps} ($\nsub(t) = 1$):
\begin{equation}
\label{eq:expansion}
\Lyap(\Wm_{t+1}) \leq \Lyap(\Wm_t) \cdot
  \left(1 - \rS \cdot \rhomincross(P_t) + \rW\right)
  + \dmax \cdot \Delta\rho(t)
\end{equation}
where $\rW = \beta(t) \cdot c_{\mathrm{D}}$ is the error-proportional
degradation rate.  On subsumption steps, the multiplicative factor may
exceed $1$ \emph{and} the additive kick may be nonzero.
\end{proposition}

\begin{proof}
The proof rests on the additive decomposition of the Lyapunov step:
\begin{equation}
\label{eq:lyap_decomposition}
\Lyap(\Wm_{t+1}) = \Lyap(\Wm_t) - \alpha(t) \cdot S(t)
  + \beta(t) \cdot D(t)
\end{equation}
where $S(t) \geq 0$ is the correction signal and $D(t) \geq 0$ is the
degradation signal.  This decomposition holds when the world-model
update separates into a correction component (driven by observation
channels) and a degradation component (driven by channel loss and
exogenous drift), with $\alpha(t)$ and $\beta(t)$ as their respective
coefficients.  The EWMA learning rules from Paper~4's risk-trust model
have exactly this structure: the update is linear in the residual (the
correction component) with a coefficient proportional to the learning
rate, and the drift component is separable by Lipschitz continuity of
$f_{\Wm}$ (Assumption~C2$'$).

For non-subsumption steps ($\nsub(t) = 0$), Assumption~C2$'$ gives
$D(t) \leq \dmax \cdot \Delta\rho(t) = 0$ (since $P_{t+1} = P_t$
when no subsumption occurs, $\Delta\rho(t) = 0$).  Applying
Assumption~C1:
\begin{align*}
\Lyap(\Wm_{t+1})
&\leq \Lyap(\Wm_t) - \alpha(t) \cdot c_{\mathrm{S}} \cdot
  \rhomincross(P_t) \cdot \Lyap(\Wm_t) \\
&= \Lyap(\Wm_t) \cdot \left(1 - \rS \cdot \rhomincross(P_t)\right).
\end{align*}

For subsumption steps ($\nsub(t) = 1$), apply both assumptions:
\begin{align*}
\Lyap(\Wm_{t+1})
&\leq \Lyap(\Wm_t)
  - \alpha(t) \cdot c_{\mathrm{S}} \cdot \rhomincross(P_t) \cdot
    \Lyap(\Wm_t)
  + \beta(t) \cdot c_{\mathrm{D}} \cdot \Lyap(\Wm_t)
  + \beta(t) \cdot \dmax \cdot \Delta\rho(t) \\
&= \Lyap(\Wm_t) \cdot
  \left(1 - \rS \cdot \rhomincross(P_t) + \rW\right)
  + \dmax \cdot \Delta\rho(t)
\end{align*}
where the last line absorbs $\beta(t)$ into $\dmax$ (redefining
$\dmax$ to include the degradation coefficient, which is permissible
since $\dmax$ is a free constant in C2$'$).
\end{proof}

\subsection{Cumulative error budget}

The per-step bound is no longer purely multiplicative (the additive
$\dmax \cdot \Delta\rho$ term prevents clean products of factors).
We therefore track the cumulative error budget over a window
$[t_0, t_0 + T]$:
\begin{equation}
\label{eq:error_budget}
B(t_0, T) =
  \underbrace{\sum_{t \in \mathrm{sub}} \dmax \cdot \Delta\rho(t)}
    _{\text{drift exposure}}
  + \underbrace{\sum_{t \in \mathrm{sub}} \rW \cdot \Lyap(\Wm_t)}
    _{\text{error-proportional degradation}}
  - \underbrace{\sum_{t=t_0}^{t_0+T-1} \rS \cdot \rhomincross(P_t) \cdot
    \Lyap(\Wm_t)}_{\text{cumulative correction}}
\end{equation}
where $\sum_{t \in \mathrm{sub}}$ ranges over the subsumption steps in
the window.

The first sum is total drift exposure from channel loss.  The second is
error-proportional degradation on subsumption steps.  The third is total
correction across all steps.  The system is net-degrading over
$[t_0, t_0 + T]$ when $B(t_0, T) > 0$.

\begin{remark}[Linearized approximation]
\label{rem:linearized}
When $\Lyap(\Wm_t)$ is approximately constant over the window and
$\Delta\rho$ is spread across $T_{\mathrm{sub}}$ subsumption steps, the
cumulative condition simplifies to
\begin{equation}
\label{eq:linearized}
\rS \cdot \rhomincross > \rW \cdot \rsub
  + \dmax \cdot \frac{\Delta\rho_{\mathrm{avg}}}{\Lyap_{\mathrm{avg}}}
\end{equation}
where $\rsub = T_{\mathrm{sub}} / T$ is the subsumption frequency and
$\Delta\rho_{\mathrm{avg}}$, $\Lyap_{\mathrm{avg}}$ are time-averaged
quantities.  The $\dmax$ term makes the condition harder to satisfy
when the rate of channel loss on well-calibrated dimensions is high
relative to current error---correctly diagnosing the failure mode where
removing a channel that was doing all the correction work appears safe
because current error is low.
\end{remark}

\begin{remark}[Why $\rhomincross$, not $\rho_{\min}$]
\label{rem:why_cross}
Proposition~\ref{prop:error_dynamics} uses cross-substrate redundancy
$\rhomincross$ rather than all-channel redundancy $\rho_{\min}$ because
the coercivity bound C1 requires correction signals to be additive
across channels.  Same-substrate channels may have correlated errors,
breaking additivity.  Cross-substrate channels satisfy Paper~5's two
distinct properties:
\begin{enumerate}
\item[(a)] \emph{Channel isolation} (unconditional): corruption on one
  substrate cannot propagate to channels on another substrate.
\item[(b)] \emph{Channel independence} (conditional on the
  joint-factorization assumption): channel outputs are statistically
  independent given the true state.
\end{enumerate}
Isolation prevents corruption propagation; independence makes correction
signals additive.  The Lyapunov bound uses additivity, which requires
independence, which requires joint factorization.  If joint factorization
fails (shared environmental noise affects both substrates), correction
signals may be sub-additive, and the effective $\rhocross_k$ should be
discounted accordingly.  Using $\rhomincross$ makes the bound valid
without assuming same-substrate independence while still requiring
cross-substrate independence via the joint-factorization assumption.
\end{remark}

\begin{remark}
The cumulative-budget structure is the core of the paper.  The phase
boundary (Section~\ref{sec:phase_boundary}) is a condition on the
\emph{frequency and severity} of subsumption relative to the
inter-shock correction capacity, not a condition that must hold at every
individual step.  Subsumption events are discrete shocks; the system can
tolerate occasional shocks as long as the inter-shock correction is
sufficient.
\end{remark}
